% =============================================================
%  Chapter 3 — Geometry-Only Structured Navigation:
%              The GRL-SNAM Foundation
%  Rewritten per professor's feedback:
%  - manifold/cotangent-bundle advantage made explicit
%  - role of p as internal evaluator/guide
%  - factored algorithm hierarchy
%  - loss terms mapped to Hamiltonian components
%  - incremental update / efficiency story
% =============================================================

\chapter{Geometry-Only Structured Navigation: The GRL-SNAM Foundation}
\label{chap:grlsnam}

\providecommand{\Ham}{\mathcal{H}}
\providecommand{\Potential}{\mathcal{R}}
\providecommand{\PhaseState}{z}
\providecommand{\Pos}{q}
\providecommand{\Mom}{p}
\providecommand{\Mass}{\mathrm{M}}
\providecommand{\J}{\mathbb{J}}

% =======================================================================
\section{Motivation: Why This Case First?}
\label{sec:ch3:motivation}
% =======================================================================

Before adding material risk or semantic structure, we must
establish what a \emph{geometry-only} Hamiltonian learner can already
do and what structural gaps geometry alone cannot close.
This chapter serves that purpose.
It introduces the four ideas that carry through this work:
the policy as a phase-space flow map, local sensing and staged
energy shaping, online cotangent-space correction, and
multi-timescale coordination.
Understanding these in the simpler geometry-only setting is necessary
before Chapter~\ref{chap:material} enriches the energy landscape with
material risk and tail-risk objectives.

% =======================================================================
\section{Geometry-Only Problem Setting}
\label{sec:ch3:setting}
% =======================================================================

\paragraph{Task.}
The agent must navigate from $\vec{q}_0$ to a goal
$\vec{x}_g \in \mathbb{R}^2$ through an unknown workspace
$\mathcal{W}$ whose occupancy field $I : \mathcal{W}\to\{0,1\}$ is
never given in full.
The robot is deformable (a ring in our experiments) with configuration
$\vec{q} = (\vec{y},\vec{c},\vec{\psi}) \in \mathcal{Q}
= \mathcal{Q}_y \times \mathcal{Q}_f \times \mathcal{Q}_o$,
where $\vec{c}\in\mathbb{R}^2$ is the pose,
$\vec{y}$ collects sensor variables, and
$\vec{\psi}$ parameterizes morphology.

\paragraph{Geometry-only restriction.}
In this chapter the environment context at stage $\tau$ contains
\emph{only geometric} information:
\begin{equation}
  \mathcal{E}_\tau
  = \bigl\{
      \vec{c}_{\mathrm{target}},\;
      C(\mathcal{E}_\tau),\;
      \{d_i(\vec{q};\mathcal{E}_\tau)\}_{i=1}^{C(\mathcal{E}_\tau)}
    \bigr\},
  \label{eq:ch3:context}
\end{equation}
where $d_i \geq 0$ are signed clearances to locally sensed obstacles.
No semantic label, material property, or latent risk signal is available.
The agent observes only the local window $\mathcal{W}(\vec{q}_t)$ of
half-width $\hat{d}$ centred at the current pose; sensing economy is
measured by the mapping ratio
\begin{equation}
  \rho_{\mathrm{map}} := \frac{\mathrm{area}\!\left(\bigcup_t \mathcal{W}(\vec{q}_t)\right)}{L^2}.
  \label{eq:ch3:map_ratio}
\end{equation}

% Teaser: Progressive Hamiltonian Landscape Learning under Local Sensing
\begin{figure}[t]
\centering
\begin{tikzpicture}[
  font=\small,
  >=Latex,
  box/.style={draw, rounded corners=2pt, very thick, inner sep=1pt, fill=black!1},
  rowtag/.style={draw, rounded corners=6pt, fill=black!6, inner xsep=6pt, inner ysep=2pt,
                 font=\scriptsize\bfseries, align=center},
  header/.style={draw, rounded corners=3pt, fill=black!10, very thick,
                 inner xsep=6pt, inner ysep=3pt, font=\scriptsize\bfseries, align=center},
  legendbox/.style={draw, rounded corners=2pt, thick, fill=white, inner sep=5pt},
  arrow/.style={->, very thick},
  accent/.style={draw=RoyalBlue, very thick},
  accent2/.style={draw=Crimson,  very thick},
  accent3/.style={draw=ForestGreen, very thick},
  alabel/.style={font=\scriptsize, fill=white, rounded corners=2pt, inner sep=2pt}
]

% ------------------- 2-column grid: Physical | Energy -------------------
\matrix (G) [matrix,
             row sep=10mm,
             column sep=38mm,
             nodes={anchor=center}] at (0,0) {

  \node[box] (phys0) {\includegraphics[width=7.0cm,height=2.35cm,keepaspectratio]{figs/phys_try1.png}}; &
  \node[box] (eng0)  {\includegraphics[width=7.0cm,height=2.35cm,keepaspectratio]{figs/energy_try1.png}}; \\

  \node[box] (phys1) {\includegraphics[width=7.0cm,height=2.35cm,keepaspectratio]{figs/phys_try2.png}}; &
  \node[box] (eng1)  {\includegraphics[width=7.0cm,height=2.35cm,keepaspectratio]{figs/energy_try2.png}}; \\

  \node[box] (phys2) {\includegraphics[width=7.0cm,height=2.35cm,keepaspectratio]{figs/phys_try3.png}}; &
  \node[box] (eng2)  {\includegraphics[width=7.0cm,height=2.35cm,keepaspectratio]{figs/energy_try3.png}}; \\
};

\node[header, text width=4.5cm, anchor=south] (colL) at ($(phys0.north)+(0,7mm)$)
{Physical space\\(unknown $\rightarrow$ revealed by local sensing)};

\node[header, text width=4.0cm, anchor=south] (colR) at ($(eng0.north)+(0,7mm)$)
{Hamiltonian energy terrain\\(learned progressively)};

% ------------------- Title -------------------
\node[
  font=\small\bfseries,
  anchor=south,
  text width=13cm,
  align=center
]
at ($(colL.north)!0.5!(colR.north)+(0,2.5mm)$)
{Progressive Hamiltonian Landscape Learning from Local Sensing\\
along Hamiltonian geodesic pathways};

% ------------------- Vertical progression arrows -------------------
\draw[arrow] (phys0.south) -- (phys1.north);
\draw[arrow] (phys1.south) -- (phys2.north);
\draw[arrow] (eng0.south) -- (eng1.north);
\draw[arrow] (eng1.south) -- (eng2.north);

% ------------------- Cross-column arrows + boxed labels -------------------
\draw[arrow, accent] (phys0.east) -- (eng0.west)
  node[midway, above=2mm, alabel] {scan local env $\Rightarrow$ build initial $H_{\tau}$};

\draw[arrow, accent2] (phys1.east) -- (eng1.west)
  node[midway, above=2mm, alabel] {Try 1 feedback $\Rightarrow$ update $H_{\tau}$};

\draw[arrow, accent3] (phys2.east) -- (eng2.west)
  node[midway, above=2mm, alabel] {Try 2 feedback $\Rightarrow$ refine $H_{\tau}$};

% ------------------- Bottom idea box -------------------
\node[draw, rounded corners=2pt, thick, fill=black!1,
      inner sep=6pt, text width=13cm, align=center, anchor=north]
  at ($(G.south)+(0,-11mm)$) {
  \textbf{Idea:} In an unknown environment, the agent repeatedly senses a \emph{local} neighborhood,
  attempts to reach the stage goal, and incrementally learns a Hamiltonian energy landscape whose valleys guide future attempts under partial observability.
};

\end{tikzpicture}
\caption{\textbf{Conceptual overview: progressive energy-landscape shaping from local sensing.}
\textit{Left}: the physical workspace is initially unknown and is revealed only through local sensing around the agent, producing a partial view of nearby obstacles and a stage goal.
\textit{Right}: from this partial context, the agent builds an initial internal “energy terrain” and then repeatedly attempts a short plan toward the stage goal; each attempt produces simple feedback (e.g., collisions/clearance and progress), which is used to refine the terrain. Over successive refinements, the terrain becomes smoother and its low-energy valleys align with traversable corridors, guiding future attempts even under partial observability.}

\label{fig:teaser_progressive_landscape}
\end{figure}


Figure~\ref{fig:teaser_progressive_landscape} shows how the energy
landscape is progressively assembled from local sensing over successive
attempts: the agent builds an initial terrain from partial context,
refines it from rollout feedback, and converges toward a landscape
whose valleys align with traversable corridors.
% =======================================================================
\section{Why Phase Space? The Manifold and Cotangent-Bundle Viewpoint}
\label{sec:ch3:phasespace}
% =======================================================================

Standard RL methods work in state-action space:
a network maps $\vec{q}_t \mapsto a_t$, the environment advances the
state, and a scalar reward drives learning.
The policy topology (which configurations attract, which repel) is
encoded implicitly in the weights with no geometric structure and
no direct handle for stable correction.

\begin{figure}[t]
    \centering
    \includegraphics[width=0.85\linewidth]{figs/flat_policy_vs_phase_space_hamiltonian.png}
    \caption{
    % \textbf{Flat policy learning versus phase-space Hamiltonian learning.}
    \textit{Left:} a black-box policy maps local observations directly to actions, with evaluation remaining external and fixed.
    \textit{Right:} GRL-SNAM represents policy as a phase-space flow, where local sensing updates an internal cotangent evaluator \(p\) that shapes the Hamiltonian vector field and smoothly guides the configuration variable \(q\). The figure highlights the shift from reactive action prediction to structured geometric dynamics with incremental correction.
    }
    \label{fig:ch3:flat_vs_phase_space}
\end{figure}

GRL-SNAM~\citep{ellendula2025grlsnam} takes a different starting point.
The configuration $\vec{q}$ moves on the configuration manifold
$\mathcal{Q}$, a smooth manifold that captures the robot's degrees of
freedom.
Attached to every point $\vec{q}$ is the cotangent space
$T^*_{\vec{q}}\mathcal{Q}$, the dual of the tangent space.
An element $p \in T^*_{\vec{q}}\mathcal{Q}$ is a \emph{covector} defined as a linear functional on tangent velocities and it represents the
generalized momentum conjugate to $\vec{q}$.

The \textbf{phase space} $T^*\mathcal{Q}$ is the union of all
$(q,p)$ pairs.
A scalar function $H : T^*\mathcal{Q} \to \mathbb{R}$ on this space
defines a Hamiltonian vector field via
\begin{equation}
  \dot{z} = \J\,\nabla_z H(z),
  \qquad
  \dot{\vec{q}} = \nabla_{\Mom} H = \Mass^{-1}\Mom,
  \qquad
  \dot{\Mom} = -\nabla_{\Pos} H,
  \label{eq:ch3:hamilton}
\end{equation}
where $z = (\vec{q},\Mom)$,
$\J = \bigl[\begin{smallmatrix}0&I\\-I&0\end{smallmatrix}\bigr]$
is the symplectic matrix acting on the \emph{full} phase gradient $\nabla_z H$
(not on $\nabla_\Mom H$ alone).
The \textbf{policy is the deterministic time-$\Delta t$ flow map} of this system:
\begin{equation}
  z_{t+1} = \Phi^{\Delta t}_H(z_t),
  \qquad
  \pi_H(z_t) := \Pi_{\mathcal{Q}}\!\left(\Phi^{\Delta t}_H(z_t)\right),
  \qquad
  z_t = (\vec{q}_t, \Mom_t).
  \label{eq:ch3:policy}
\end{equation}

\paragraph{Why this matters computationally.}
In flat state-action RL, adapting the policy means retraining
network weights which is a global, potentially destructive update.
In the Hamiltonian viewpoint, adapting the policy means reshaping
$H$: adjusting the coefficients $(\beta, \{\alpha_i\}, \gamma)$
that modulate its energy terms.
Each coefficient controls a specific force channel.
The force field topology i.e.; goal-attracting basins, obstacle-repulsive
walls is explicit, inspectable, and surgically correctable.

A second advantage is structure-preserving integration.
For the conservative Hamiltonian substep (without dissipation or port
input), the leapfrog integrator is a \emph{symplectomorphism}: it
preserves the discrete symplectic form $\omega = \sum_i dq^i\wedge dp_i$
up to roundoff, controlling long-run phase-space volume drift that
accumulates in generic Euler-step integrators~\citep{reich1996symplectic,hairer2006geometric}.
The dissipative port ($R_\theta$) and correction input ($u^\xi$) are
applied as separate substeps; the full port-Hamiltonian update is not
symplectic but satisfies the energy-dissipation inequality
$\dot{H} \leq y^\top u$ (passivity).

The third advantage is that the value function is related to the
costate: when the Hamilton–Jacobi–Bellman equation admits a smooth
solution $\mathcal{V}^*(t,\vec{q};\mathcal{E})$,
the costate satisfies
$\Mom(t) = \nabla_{\Pos}\mathcal{V}^*(t,\vec{q}(t);\mathcal{E})$.
Conventional RL approximates $\mathcal{V}^*$ via Bellman backups which is a
backward-time PDE solved approximately.
GRL-SNAM instead learns the reduced Hamiltonian directly, producing
forward-time rollouts without Bellman recursion.

% =======================================================================
\section{The Role of \texorpdfstring{$p$}{p}: Internal Evaluator,
         Predictor, and Guide}
\label{sec:ch3:p_role}
% =======================================================================

The momentum variable $\Mom$ is often introduced as "the conjugate
variable to $\Pos$" and left at that.
This undersells its computational role.
In GRL-SNAM, $\Mom$ simultaneously performs five functions at each
integration step.
\begin{figure}[htbp]
    \centering
    \includegraphics[width=\linewidth]{figs/role_of_p_one_update.png}
    \caption{
    \textbf{The role of \(p\) in one phase-space update.}
    Local sensing feeds the cotangent variable \(p_t\), which serves as an internal evaluator and force predictor.
    Through the leapfrog half-step \(p_{t+\frac12}\), this information is converted into a smooth update of the configuration variable \(q\), yielding the corrected transition \((q_t,p_t)\mapsto(q_{t+1},p_{t+1})\).
    }
    \label{fig:ch3:role_of_p_update}
\end{figure}
\paragraph{1. Geometric accumulator.}
$\Mom_t$ encodes a summary of the geometry sensed so far,
weighted by the current barrier weights $\{\alpha_i^t\}$.
Through \eqref{eq:ch3:hamilton}, each barrier clearance $d_i$
pushes the costate via $-\nabla_{\Pos}b_i(d_i)$;
a long sequence of clearance signals accumulates as a running
cotangent-space summary of local obstacle geometry.

\paragraph{2. Progress evaluator.}
The inner product $\Mom_t^\top \Mass^{-1}\Mom_t = 2H_{\mathrm{kin}}$
is the current kinetic energy.
When $H_{\mathrm{kin}}$ is large but goal distance is not decreasing,
the agent is spending kinetic energy unproductively which is an early
warning of oscillation or stall.
This is the information used in the passivity surrogate
$\mathcal{L}_{\mathrm{pass}} = \frac{1}{T}\sum_t\mathrm{softplus}(\Delta H_t)$.

\paragraph{3. Force predictor.}
From \eqref{eq:ch3:hamilton}, $\dot{\vec{q}} = \Mass^{-1}\Mom_t$:
the current momentum \emph{predicts} the next configuration
displacement before that displacement is computed.
This half-step predictive role is exactly what the leapfrog
integrator exploits: update $\Mom$ using the current gradient,
then advance $\Pos$ using the updated $\Mom$.
The costate is always one half-step ahead.

\paragraph{4. Online correction receiver.}
The online adaptation law updates $\{\alpha_i^t\}$ from the observed
residual $\Delta y^{\mathrm{des}}_t - J_t\Delta\zeta_t$.
This residual enters the dynamics through the costate equation
$\dot\Mom = -\nabla_{\Pos}\Potential - R_\theta$:
the correction reshapes the force field by changing the gradient
that acts on $\Mom$, not by directly perturbing $\Pos$.
The costate is the \emph{receiver} of adaptation signals.

\paragraph{5. Port correction amender.}
When energy-coefficient updates alone cannot satisfy the desired
observable change (large residual $r_t$), a port correction
$u_f^\xi$ is injected directly into the costate equation.
This is the non-conservative port of the port-Hamiltonian structure:
$\dot\Mom \leftarrow \dot\Mom + G_f^\xi u_f^\xi$.
The evaluator in conventional RL is external (a separate value
network).
In GRL-SNAM, the evaluator is the costate $\Mom$: it lives
\emph{inside} the evolving trajectory.

\begin{quote}\itshape
$\vec{q}$ is the primal movement variable, advancing on the
configuration manifold.
$\Mom$ is the cotangent-space intelligence: it senses geometry,
evaluates progress, predicts the next force balance, receives
adaptation signals, and corrects the motion field from within
the rollout and not from an external critic.
\end{quote}

% =======================================================================
\section{Reduced Hamiltonian and Energy Decomposition}
\label{sec:ch3:hamiltonian}
% =======================================================================

\subsection{From OCP to Reduced Hamiltonian}
\label{subsec:ch3:ocp}

At stage $\tau$, navigation is a constrained OCP
(see §\ref{subsec:prelim:pmp} for the full derivation).
Following the mechanical surrogate construction of
§\ref{subsec:prelim:pmp}, we adopt a
\textbf{structured mechanical Hamiltonian} whose kinetic--potential
separation is motivated by the PMP costate structure but learned
end-to-end rather than derived as an exact PMP consequence:
\begin{equation}
  H(\vec{q},\Mom;\mathcal{E}_\tau)
  = \underbrace{\tfrac{1}{2}\Mom^\top\Mass^{-1}\Mom}_{H_{\mathrm{kin}}}
  + \underbrace{\Potential(\vec{q};\mathcal{E}_\tau)}_{H_{\mathrm{pot}}},
  \label{eq:ch3:ham}
\end{equation}
where the kinetic term drives forward momentum and the potential
$\Potential$ encodes navigation objectives.
The optimal trajectory follows Hamilton's equations~\eqref{eq:ch3:hamilton}.

\subsection{The Geometry-Only Potential: Term-by-Term}
\label{subsec:ch3:potential}

The potential $\Potential$ decomposes additively over four terms.
Each term corresponds to one navigation objective; each term contributes
one force channel; each force channel has one learnable coefficient.

\begin{equation}
  \Potential(\vec{q};\mathcal{E}_\tau)
  = \underbrace{E_{\mathrm{sensor}}(\vec{y})}_{\text{(i) sensing}}
  + \underbrace{\beta^t E_{\mathrm{goal}}(\vec{c},\vec{c}_{\mathrm{target}})}_{\text{(ii) goal}}
  + \underbrace{\lambda^t E_{\mathrm{obj}}(\vec{\psi})}_{\text{(iii) shape}}
  + \underbrace{\sum_{i\in\mathcal{C}_t} \alpha_i^t\,b(d_i(\vec{q};\mathcal{E}_\tau))}_{\text{(iv) barrier}},
  \label{eq:ch3:potential}
\end{equation}
where $b(\cdot)$ is an IPC-style barrier potential and the active
constraint set $\mathcal{C}_t = \{i \mid d_i \leq \hat{d}\}$ is
discovered online by sensing.

\begin{table}[htbp]
\centering
\caption{%
  \textbf{Energy terms, their force channels, coefficients, and
  what each shapes in the dynamics.}
  Each row is one link in the enrichment chain
  $H_i \to F_i \to \vartheta_i$.
}
\label{tab:ch3:terms}
\small
\setlength{\tabcolsep}{5pt}
\renewcommand{\arraystretch}{1.3}
\resizebox{\textwidth}{!}{
\begin{tabular}{lllll}
\toprule
\textbf{Term} & \textbf{Force $F_i = -\nabla_{\Pos}H_i$} &
\textbf{Coeff.} & \textbf{Shapes} & \textbf{Updated by} \\
\midrule
$E_{\mathrm{sensor}}$ & $-\nabla_{\vec{y}}\|\vec{y}\|_S^2$
  & $\omega_y$ (fixed) & Sensor alignment & Offline only \\
$E_{\mathrm{goal}}$ & $-2\beta(\vec{c}-\vec{c}_{\mathrm{target}})$
  & $\beta^t$ & Goal-attracting basin & $\mathcal{L}_{\mathrm{traj}}$ \\
$E_{\mathrm{obj}}$ & $-\lambda\nabla_{\vec\psi}E_{\mathrm{obj}}$
  & $\lambda^t$ & Shape/deformation & $\mathcal{L}_{\mathrm{traj}}$ \\
$b(d_i)$ & $-\alpha_i\nabla_{\Pos}b(d_i)$
  & $\{\alpha_i^t\}$ & Repulsive walls near obstacles & Online $J_t$ update \\
$R_\theta = \gamma\Mass^{-1}\Mom$ & (dissipation, not $\nabla H$)
  & $\gamma$ & Energy removal, passivity & $\mathcal{L}_{\mathrm{fric}}$ \\
\bottomrule
\end{tabular}}
\end{table}

Table~\ref{tab:ch3:terms} summarizes the five channels.
Note that the dissipative port $R_\theta = \gamma\Mass^{-1}\Mom$
is \emph{not} part of $H$, it enters the costate equation
separately as $\dot\Mom \leftarrow \dot\Mom - \tau R_\theta$,
preserving the port-Hamiltonian separation of stored energy and
dissipation (Definition~\ref{def:passivity}).

\subsection{The Meta-Navigator: Context to Coefficients}
\label{subsec:ch3:metanav}

The meta-navigator $g_\xi$ maps each local context to a full
coefficient vector in a single forward pass:
\begin{equation}
  g_\xi(\mathcal{E}_\tau)
  = \bigl[\,\eta_\xi(\mathcal{E}_\tau),\;\mu_\xi(\mathcal{E}_\tau),\;u_f^\xi\,\bigr],
  \qquad
  \eta_\xi = (\beta, \lambda, \{\alpha_i\}_{i\in\mathcal{C}_t}).
  \label{eq:ch3:metanav}
\end{equation}
The parameters $\xi$ are learned offline; $g_\xi$ runs as a
fixed network at inference time where no optimization occurs per query.
This is the first computational efficiency property: \emph{the cost
per stage is one forward pass through $g_\xi$, not a per-stage
optimization}.

The output $\eta_\xi$ provides the inter-term tradeoff weights;
$\mu_\xi$ provides the friction coefficient for the dissipative port;
$u_f^\xi$ provides a port correction when energy reweighting alone
is insufficient (see §\ref{sec:ch3:online}).

\begin{figure}[t]
\centering
\resizebox{\linewidth}{!}{%
\begin{tikzpicture}[
  node distance=0.5cm,
  font=\small,
  % --- Colors ---
  sensorcolor/.style={draw=RoyalBlue, fill=RoyalBlue!10},
  framecolor/.style ={draw=Crimson,   fill=Crimson!10},
  shapecolor/.style ={draw=ForestGreen,fill=ForestGreen!10},
  navcolor/.style   ={draw=Orange,    fill=Orange!10},
  band/.style       ={fill=black!3, draw=none},
  % --- Styles ---
  mod/.style   ={rectangle,rounded corners=2pt,draw,very thick,align=center,
                 minimum width=4.8cm,minimum height=2.2cm,#1},
  nav/.style   ={rectangle,rounded corners=2pt,draw,very thick,align=center,
                 minimum width=14.0cm,minimum height=2.0cm,#1},
  io/.style    ={circle,draw,thick,inner sep=0pt,minimum size=2.8mm,fill=white},
  qarr/.style  ={thick,dashed,-{Stealth[length=2.2mm]}},
  rarr/.style  ={thick,-{Stealth[length=2.2mm]}},
  chip/.style  ={rectangle,rounded corners=1pt,draw,thin,inner sep=2pt,
                 font=\scriptsize,fill=white},
  dsicon/.style={rectangle,draw,thin,fill=white,minimum width=4.5mm,minimum height=4.5mm},
  badge/.style ={circle,draw,thin,fill=white,minimum size=5mm,font=\scriptsize\bfseries},
  xtiny/.style ={font=\scriptsize},
  response/.style={rectangle,draw,thin,fill=white,rounded corners=1pt,font=\scriptsize,align=left,inner sep=3pt}
]

% ====== GRID COORDS ======
\coordinate (ColY) at (-5.2,0);   % Sensor column x
\coordinate (ColF) at ( 0.0,0);   % Frame column x
\coordinate (ColO) at ( 5.2,0);   % Shape column x

% ====== Background band ======
\node[rectangle,band,minimum width=16.5cm,minimum height=2.8cm,anchor=center] at (0,6.0) {};

% ====== MODULES (moved up) ======
% Sensor module
\node[mod=sensorcolor,anchor=center] (sensor) at ($(ColY)+(0,6.0)$) {
    {\bfseries\color{RoyalBlue}Sensor Policy $\pi_y^{\theta_y,\xi}$}\\[2pt]
    {\scriptsize $s_y^{\theta_y,\xi}(z_y,t;\mathcal{E})=\mathbb{J}\nabla_{z_y}h_y^{\theta_y}$}\\[1pt]
    {\scriptsize\color{gray} $\Gamma_y^\xi\equiv 0,\; G_y^\xi u_y^\xi\equiv 0$}
};

% Frame module
\node[mod=framecolor,anchor=center] (frame) at ($(ColF)+(0,6.0)$) {
    {\bfseries\color{Crimson}Frame Policy $\pi_f^{\theta_f,\xi}$}\\[1pt]
    {\scriptsize $s_f^{\theta_f,\xi}=\mathbb{J}\nabla_{z_f}h_f^{\theta_f}+\begin{bmatrix}0\\-\Gamma_f^\xi\dot{q}_f+G_f^\xi u_f^\xi\end{bmatrix}$}\\[1pt]
    {\scriptsize\color{gray} $\Gamma_f^\xi=\mu_\xi(\mathcal{E})\mathbf{I},\; G_f^\xi=\mathbf{I}$}
};

% Shape module
\node[mod=shapecolor,anchor=center] (shape) at ($(ColO)+(0,6.0)$) {
    {\bfseries\color{ForestGreen}Shape Policy $\pi_o^{\theta_o,\xi}$}\\[2pt]
    {\scriptsize $s_o^{\theta_o,\xi}(z_o,t;\mathcal{E})=\mathbb{J}\nabla_{z_o}h_o^{\theta_o}$}\\[1pt]
    {\scriptsize\color{gray} $\Gamma_o^\xi\equiv 0,\; G_o^\xi u_o^\xi\equiv 0$}
};

% ====== Configuration space icons (top-left inside each) ======
\foreach \M/\L in {sensor/$\mathcal{Q}_y$,frame/$\mathcal{Q}_f$,shape/$\mathcal{Q}_o$}{
  \node[dsicon] (ds-\M) at ([xshift=4mm,yshift=-3mm]\M.north west) {};
  \draw[thin] (ds-\M.south west)++(0,0.6mm) rectangle ++(4.5mm,0.8mm);
  \draw[thin] (ds-\M.south west)++(0,1.8mm) rectangle ++(4.5mm,0.8mm);
  \draw[thin] (ds-\M.south west)++(0,3.0mm) rectangle ++(4.5mm,0.8mm);
  \node[xtiny] at ([xshift=7mm]ds-\M.east) {\L};
}

% ====== Time badges (outside top-right) ======
\node[badge,text=RoyalBlue]   at ([xshift=-14pt,yshift=8pt]sensor.north east) {$T_y$};
\node[badge,text=Crimson]     at ([xshift=-14pt,yshift=8pt]frame.north east)  {$T_f$};
\node[badge,text=ForestGreen] at ([xshift=-14pt,yshift=8pt]shape.north east)  {$T_o$};

% ====== Timescale annotation ======
\node[xtiny,anchor=west] at (8.0,7.6) {$T_y \gg T_f \gg T_o$};

% ====== POTENTIAL TERMS (above modules) ======
\node[chip,fill=RoyalBlue!5] at ([yshift=6pt]sensor.north) {$\mathcal{R}_y: E_{\mathrm{sensor}} + \sum_i\alpha_i b(d_i)$};
\node[chip,fill=Crimson!5] at ([yshift=6pt]frame.north) {$\mathcal{R}_f: \beta E_{\mathrm{goal}} + \sum_i\alpha_i b(d_i)$};
\node[chip,fill=ForestGreen!5] at ([yshift=6pt]shape.north) {$\mathcal{R}_o: \lambda E_{\mathrm{obj}} + \sum_i\alpha_i b(d_i)$};

% ====== RESPONSE BOXES (between modules and navigator) ======
\node[response,fill=RoyalBlue!5] (Ry) at ($(ColY)+(0,3.6)$) {
    $\mathsf{R}_y=\big\{z_{0:T_y}^{(y)},\mathrm{QoI}_y\big\}$\\
    $\mathrm{QoI}_y=\Delta\mathcal{E}$
};
\node[response,fill=Crimson!5] (Rf) at ($(ColF)+(0,3.6)$) {
    $\mathsf{R}_f=\big\{z_{0:T_f}^{(f)},\mathrm{QoI}_f\big\}$\\
    $\mathrm{QoI}_f=\{v_t^{(f)}\}_{t=0}^{T_f}$
};
\node[response,fill=ForestGreen!5] (Ro) at ($(ColO)+(0,3.6)$) {
    $\mathsf{R}_o=\big\{z_{0:T_o}^{(o)},\mathrm{QoI}_o\big\}$\\
    $\mathrm{QoI}_o=\{\min_i d_i\}_{t=0}^{T_o}$
};

% ====== NAVIGATOR ======
\node[nav=navcolor,anchor=center] (nav) at ($(0,1.2)$) {
  {\bfseries Navigator $g_{\xi}$ (Meta-Hamiltonian Composer)}\\[2pt]
  {\scriptsize Tokenize: $\mathcal{T}_t=\mathrm{Tokenize}(\mathcal{C}_t(\mathcal{E},q_t),\,q_t,\,\mathbf{x}_g)$\quad\quad Output: $g_\xi(\mathcal{E})=\big[\,\eta_\xi,\;\mu_\xi,\;u_f^\xi\,\big]$}\\[1pt]
  {\scriptsize Dual weights: $\eta_\xi(\mathcal{E})=(\beta,\lambda,\{\alpha_i\}_{i\in\mathcal{C}_t})$\quad\quad Surrogate: $H=\tfrac{1}{2}p^{\top}M^{-1}p+\mathcal{R}(q;\eta_\xi,\mathcal{E})$}
};

% ====== QUERY ARROWS (Navigator to Modules) ======
\draw[qarr,draw=RoyalBlue]
  ([xshift=-8pt]nav.north-|ColY) -- node[xtiny,fill=white,inner sep=1pt,pos=0.5,left]
  {$(H_y,z_y^t,T_y)$} ([xshift=-8pt]Ry.south-|ColY);

\draw[qarr,draw=Crimson]
  ([xshift=-8pt]nav.north-|ColF) -- node[xtiny,fill=white,inner sep=1pt,pos=0.5,left]
  {$(H_f,z_f^t,T_f,\Gamma_f^\xi,u_f^\xi)$} ([xshift=-8pt]Rf.south-|ColF);

\draw[qarr,draw=ForestGreen]
  ([xshift=-8pt]nav.north-|ColO) -- node[xtiny,fill=white,inner sep=1pt,pos=0.5,left]
  {$(H_o,z_o^t,T_o)$} ([xshift=-8pt]Ro.south-|ColO);

% ====== RESPONSE ARROWS (Response boxes to Navigator) ======
\draw[rarr,draw=RoyalBlue]
  ([xshift=8pt]Ry.south-|ColY) -- ([xshift=8pt]nav.north-|ColY);

\draw[rarr,draw=Crimson]
  ([xshift=8pt]Rf.south-|ColF) -- ([xshift=8pt]nav.north-|ColF);

\draw[rarr,draw=ForestGreen]
  ([xshift=8pt]Ro.south-|ColO) -- ([xshift=8pt]nav.north-|ColO);

% ====== ARROWS from Modules to Response boxes ======
\draw[rarr,draw=RoyalBlue] (sensor.south) -- (Ry.north);
\draw[rarr,draw=Crimson] (frame.south) -- (Rf.north);
\draw[rarr,draw=ForestGreen] (shape.south) -- (Ro.north);

% ====== Environment & Trajectory ======
\node[rectangle,draw,thick,fill=white,minimum width=2.6cm,minimum height=1.0cm,align=center]
      (env) at (-9.2,1.2) {{\bfseries Environment $\mathcal{E}$}\\ {\scriptsize $z_0,\;\mathbf{x}_g,\;\{d_i\geq 0\}$}};
\node[rectangle,draw,thick,fill=white,minimum width=2.6cm,minimum height=1.0cm,align=center]
      (traj) at ( 9.2,1.2) {{\bfseries Trajectory}\\ {\scriptsize $\{z_t\}_{t=0}^{T}$}};

\draw[thick,-{Stealth[length=2.5mm]}] (env.east) -- (nav.west);
\draw[thick,-{Stealth[length=2.5mm]}] (nav.east) -- (traj.west);

% ====== Side annotations ======
% Symplectic matrix definition (left)
\node[rectangle,draw,thin,fill=white,align=left,font=\scriptsize,inner sep=4pt]
      at (-9.8,5.8) {
       \textbf{Symplectic form:}\\[1pt]
       $\mathbb{J}=\begin{bmatrix}0 & I\\-I & 0\end{bmatrix}$\\[1pt]
       $\mathbb{J}\nabla_z h = \begin{bmatrix}\nabla_p h\\-\nabla_q h\end{bmatrix}$
      };

% Online adaptation (right)
\node[rectangle,draw,thin,fill=white,align=left,font=\scriptsize,inner sep=4pt]
      at (10.5,5.8) {
       \textbf{Online observables:}\\[1pt]
       $y_t=[-\mathrm{clr}_t,\mathrm{dist}_t,-\mathrm{speed}_t]^\top$\\[1pt]
       Adapt $\zeta_t=[\eta_t,\mu_t]$ via $J_t\approx\partial y/\partial\zeta$
      };

% ====== Compact Legend ======
\node[rectangle,draw,thick,fill=white,align=left,font=\scriptsize,inner sep=4pt]
      at (0,-0.8) {
       \textbf{Legend:}\;
       \tikz[baseline=-0.5ex]{\node[io]{};} port\quad
       \tikz[baseline=-0.5ex]{\draw[qarr] (0,0)--(8mm,0);} query\quad
       \tikz[baseline=-0.5ex]{\draw[rarr] (0,0)--(8mm,0);} response
       \qquad
       \textbf{Policy as flow:}\; $\pi_k^{\theta_k,\xi}(\mathcal{E}):=\Phi_k^{\Delta t}(\cdot;\mathcal{E})$
      };

\end{tikzpicture}}
\caption{%
Modular score-field architecture for GRL-SNAM.
The Navigator $g_\xi$ tokenizes active constraints $\mathcal{C}_t$, state $q_t$, and goal $\mathbf{x}_g$, then outputs dual weights $\eta_\xi=(\beta,\lambda,\{\alpha_i\})$, friction $\mu_\xi$, and port correction $u_f^\xi$.
Each module $k\in\{y,f,o\}$ integrates its score field $s_k^{\theta_k,\xi}=\mathbb{J}\nabla_{z_k}h_k^{\theta_k}$ (plus non-conservative terms for the frame module) over horizon $T_k$ and returns response $\mathsf{R}_k$ with phase-space rollouts and QoIs.
The frame module uniquely receives dissipation $\Gamma_f^\xi$ and port input $u_f^\xi$; sensor and shape modules have purely symplectic dynamics.
}
\label{fig:independent-score-architecture}
\end{figure}

\subsection{Online Adaptation via Cotangent Feedback}
\label{sec:ch3:online}
 
§\ref{sec:ch3:p_role} established that $\Mom_t$ carries
evaluative information about the current rollout: kinetic energy,
clearance-gradient accumulation, and speed.
Online adaptation closes the loop, it reads that information
back out of $\Mom_t$ and uses it to correct the active coefficients
$\zeta_t = (\beta_t,\lambda_t,\{\alpha_i\}_{i\in\mathcal{C}_t},\mu_t)$
without touching $\xi$ or rerunning $g_\xi$.

\begin{figure*}[htbp]
\centering

% -------- Left subfigure --------
\begin{subfigure}[htbp]{0.8\textwidth}
\centering
\setlength{\tabcolsep}{2pt}
\small\sffamily
\begin{tabular}{@{}c ccc@{}}
  & \textbf{A} & \textbf{B} & \textbf{C} \\[4pt]
  
  \rotatebox{90}{\parbox{18mm}{\centering\textcolor{red!65!black}{\textbf{Baseline}}}} &
  \includegraphics[width=0.28\linewidth]{baseline_A.png} &
  \includegraphics[width=0.28\linewidth]{baseline_B.png} &
  \includegraphics[width=0.28\linewidth]{baseline_D.png} \\[4pt]
  
  \rotatebox{90}{\parbox{18mm}{\centering\textcolor{green!50!black}{\textbf{Ours}}}} &
  \includegraphics[width=0.28\linewidth]{ours_A.png} &
  \includegraphics[width=0.28\linewidth]{ours_B.png} &
  \includegraphics[width=0.28\linewidth]{ours_D.png} \\
\end{tabular}

\vspace{1mm}
{\scriptsize
\textcolor{red!65!black}{$\times$~fails (collision/loop)}
\hfill
\textcolor{green!50!black}{$\checkmark$~reaches goal}
}

\caption{Fixed policy vs.\ online adaptation. Baseline gets trapped or collides; our method adapts online and continues toward the goal.}
\end{subfigure}
\hfill
% -------- Right subfigure --------
\begin{subfigure}[htbp]{0.75\textwidth}
\centering
\includegraphics[width=\linewidth]{figs/teaser_online_correction_mechanism.png}
\caption{Online correction. Failures trigger local updates using observable signals (clearance, distance, speed), adjusting weights to restore safety while preserving progress.}
\end{subfigure}

\caption{\textbf{Online adaptation for mapless navigation.}}
\label{fig:teaser_combined}
\end{figure*}

The three evaluative signals already in the phase-space state become
the observable vector:
\begin{equation}
  y_t
  = \bigl[-\,\mathrm{clr}_t,\;\mathrm{dist}_t,\;-\|\Mass^{-1}\Mom_t\|\bigr]^\top,
  \qquad
  \Delta y_t^{\mathrm{des}} = y_t^\star - y_t,
  \label{eq:ch3:obs}
\end{equation}
where $\mathrm{clr}_t = \min_{i\in\mathcal{C}_t}d_i$ and
$\mathrm{dist}_t = |\vec{c}_t - \vec{x}_g|$.
No second query to $g_\xi$ is needed; the signal is already in the
trajectory.
 
To move $y_t$ toward $y_t^\star$, the Jacobian
$J_t = \partial y_t/\partial\zeta_t$ is needed.
Rather than differentiating through the integrator, it is accumulated
as a rank-1 secant update from consecutive $(y,\zeta)$ pairs:
\begin{equation}
  J_t = \rho\,J_{t-1}
        + (1-\rho)\,
          \frac{(y_t-y_{t-1})(\zeta_t-\zeta_{t-1})^\top}
               {\|\zeta_t-\zeta_{t-1}\|^2+\varepsilon}.
  \label{eq:ch3:secant}
\end{equation}
With $J_t$ in hand, the coefficient update is a
Tikhonov-regularized step:
\begin{equation}
  \Delta\zeta_t = (J_t^\top J_t + \lambda_J I)^{-1}J_t^\top\Delta y_t^{\mathrm{des}},
  \qquad
  \zeta_{t+1} = \Pi_{\mathbb{R}_+}(\zeta_t + \kappa\,\Delta\zeta_t).
  \label{eq:ch3:tikhonov}
\end{equation}
The projection $\Pi_{\mathbb{R}_+}$ preserves non-negativity of all
barrier weights, maintaining the repulsive character of each $\alpha_i$.
The solve is $(2+|\mathcal{C}_t|)$-dimensional (typically 4--10)
so the cost is $O(|\mathcal{C}_t|^3)$, not $O(|\xi|^3)$.
$J_t$ itself requires only a rank-1 update
($O(3|\mathcal{C}_t|)$) and is cached between steps.
 
When the residual $r_t = \Delta y_t^{\mathrm{des}} - J_t\Delta\zeta_t$
is non-negligible meaning the desired observable change lies
partly outside the range of $J_t$, a port impulse amends it
directly through the costate equation:
\begin{equation}
  u_{f,t}^\xi = (P_t^\top P_t + \lambda_u I)^{-1}P_t^\top r_t,
  \label{eq:ch3:port}
\end{equation}
injected as $+\tfrac{\Delta t}{2}u_{f,t}^\xi$ into the half-step
momentum \eqref{eq:ch3:half_mom}.
This leaves $H$ unchanged as it is a non-conservative correction
through the port, not a modification of the stored energy.
 
The chain is therefore: $\Mom_t$ carries evaluative state
$\to$ observables are extracted $\to$ active coefficients are
updated via the secant Jacobian $\to$ residual mismatch is
absorbed by the port.
Each step reads from or writes to the evolving costate; no
external critic or second network is involved.

% =======================================================================
\section{Algorithm}
\label{sec:ch3:algorithm}
% =======================================================================

The GRL-SNAM computation is organized as two algorithms plus three
sub-procedures.
This factored structure makes explicit \emph{what runs where, what is updated, and what is cached}. 

\subsection{Top-Level Control Loop}
\label{subsec:ch3:toplevel}

Algorithm~\ref{algo:ch3:toplevel} is the outermost loop.
Its only job is to advance stages, maintain the goal check, and
delegate to the per-stage procedure.
It has no internal optimization and caches only the current
phase-space state $\PhaseState_\tau$ and context $\mathcal{E}_\tau$
across stages.

\begin{algorithm}[t]
\caption{GRL-SNAM: Top-Level Navigation Loop}
\label{algo:ch3:toplevel}
\begin{algorithmic}[1]
\State \textbf{Input:} $\vec{x}_g$, initial $\PhaseState_0 = (\vec{q}_0,\mathbf{0})$,
       trained $g_\xi$, max stages $S_{\max}$
\State $\tau \leftarrow 0$;\; $\mathcal{E}_0 \leftarrow$ \textsc{Sense}$(\vec{q}_0)$
\While{$|\vec{c}_\tau - \vec{x}_g| > \epsilon_g$ \textbf{and} $\tau < S_{\max}$}
  \State $[\eta_\xi, \mu_\xi, u_f^\xi] \leftarrow g_\xi(\mathcal{E}_\tau)$
         \hfill $\triangleright$ \emph{one forward pass; no optimization}
  \State Assemble $H(\cdot;\mathcal{E}_\tau)$ from $\eta_\xi$
         via \eqref{eq:ch3:ham}--\eqref{eq:ch3:potential}
  \State $\PhaseState_{\tau+1}, \mathcal{E}_{\tau+1}
         \leftarrow$ \textsc{StageRollout}$(\PhaseState_\tau, \mathcal{E}_\tau,
         H, \mu_\xi, u_f^\xi)$
         \hfill $\triangleright$ Algorithm~\ref{algo:ch3:stage}
  \State $\tau \leftarrow \tau + 1$
\EndWhile
\State \textbf{Return} trajectory $\{\PhaseState_t\}$, parameter history
\end{algorithmic}
\end{algorithm}

\subsection{Per-Stage Sub-Procedure}
\label{subsec:ch3:stage}

Algorithm~\ref{algo:ch3:stage} handles one stage: it integrates
the phase-space dynamics, runs the online coefficient update at
each step, and returns the updated context.

\begin{algorithm}[t]
\caption{GRL-SNAM: Per-Stage Rollout and Online Update}
\label{algo:ch3:stage}
\begin{algorithmic}[1]
\Require $\PhaseState_0^{(\tau)}, \mathcal{E}_\tau, H, \mu_\xi, u_f^\xi$,
         horizon $(T_y, T_f, T_o)$, step $\Delta t$,
         smoothed Jacobian $J_{\tau-1}$ (cached from last stage)
\State $n \leftarrow 0$;\; $\mathcal{C}_0 \leftarrow \{i : d_i(\vec{q}_0;\mathcal{E}_\tau) \leq \hat{d}\}$
\While{$n < T_f$ \textbf{and} $|\vec{c}_n - \vec{c}_{\mathrm{target}}| > \epsilon_s$}
  \If{$n \bmod T_y = 0$}
    \State \textsc{SenseUpdate}$(\vec{q}_n) \to \mathcal{E}_n$,
           update $\mathcal{C}_n$, assemble $H_y$
           \hfill $\triangleright$ Box~\ref{box:sense}
  \EndIf
  \State \textsc{SymplecticStep}$(\PhaseState_n, H, \mu_\xi, \Delta t)
         \to \PhaseState_{n+1}$
         \hfill $\triangleright$ Box~\ref{box:symp}
  \State $y_n \leftarrow [-\mathrm{clr}_n,\;\mathrm{dist}_n,\;-\mathrm{speed}_n]^\top$
  \State \textsc{OnlineCoeffUpdate}$(\PhaseState_n, y_n, J_{n-1})
         \to (\zeta_{n+1}, J_n, u_{f,n}^\xi)$
         \hfill $\triangleright$ Box~\ref{box:online}
  \State $n \leftarrow n + 1$
\EndWhile
\State \textbf{Return} $\PhaseState_n$, $\mathcal{E}_n$, $J_n$ (cache for next stage)
\end{algorithmic}
\end{algorithm}

The three sub-procedures are given as named boxes to separate
conceptual roles without adding heavyweight algorithm environments.

\begin{mdframed}[linecolor=teal!60!black, linewidth=1pt,
                 frametitle={\small\sffamily\bfseries
                 Box 3.6.1: \textsc{SenseUpdate}},
                 frametitlefont=\small\sffamily\bfseries]
\label{box:sense}
\textbf{Input:} current pose $\vec{q}_n$\\
\textbf{Cached:} previous context $\mathcal{E}_{n-1}$\\
\textbf{Computes:}
\begin{enumerate}[(i)]
  \item Query sensor: returns obstacle list
        $\{(\vec{c}_i,r_i)\}_{i=1}^{C}$ within window $\mathcal{W}(\vec{q}_n)$.
  \item Compute signed clearances
        $d_i(\vec{q};\mathcal{E}_n) = \|\vec{q}_f - \vec{c}_i\| - r_i$.
  \item Update active set $\mathcal{C}_n \leftarrow \{i : d_i \leq \hat{d}\}$.
  \item Update local goal $\vec{c}_{\mathrm{target}}(\mathcal{E}_n)$
        if stage exit reached.
\end{enumerate}
\textbf{Output:} updated $\mathcal{E}_n$, $\mathcal{C}_n$\\
\textbf{Cached after:} $\mathcal{E}_n$ replaces $\mathcal{E}_{n-1}$;
$\{\alpha_i : i \notin \mathcal{C}_n\}$ are zeroed (not stored).\\
\textbf{Cost:} $O(C(\mathcal{E}_n))$ distance evaluations.
\end{mdframed}

\vspace{6pt}

\begin{mdframed}[linecolor=orange!70!black, linewidth=1pt,
                 frametitle={\small\sffamily\bfseries
                 Box 3.6.2: \textsc{SymplecticStep}}]
\label{box:symp}
\textbf{Input:} $(\vec{q}_n, \Mom_n)$, Hamiltonian $H$,
dissipation $\mu_\xi$, step $\Delta t$\\
\textbf{Leapfrog integration} (symplecticity-preserving):
\begin{align}
  \Mom_{n+\frac{1}{2}}
  &= \Mom_n
     - \tfrac{\Delta t}{2}\nabla_{\Pos}H(\vec{q}_n,\Mom_n)
     - \tfrac{\Delta t}{2}\,\mu_\xi\Mass^{-1}\Mom_n
     + \tfrac{\Delta t}{2}\,u_{f,n}^\xi,
  \label{eq:ch3:half_mom} \\
  \vec{q}_{n+1}
  &= \vec{q}_n + \Delta t\,\Mass^{-1}\Mom_{n+\frac{1}{2}},
  \label{eq:ch3:q_update} \\
  \Mom_{n+1}
  &= \Mom_{n+\frac{1}{2}}
     - \tfrac{\Delta t}{2}\nabla_{\Pos}H(\vec{q}_{n+1},\Mom_{n+\frac{1}{2}})
     - \tfrac{\Delta t}{2}\,\mu_\xi\Mass^{-1}\Mom_{n+\frac{1}{2}}.
  \label{eq:ch3:full_mom}
\end{align}
The half-step momentum $\Mom_{n+\frac12}$ is the \emph{predictor}:
it uses the current force balance to anticipate the next position
before committing to $\vec{q}_{n+1}$.
This is the computational realization of $p$'s role as
\emph{force predictor} (§\ref{sec:ch3:p_role}).\\
\textbf{Output:} $(\vec{q}_{n+1}, \Mom_{n+1})$\\
\textbf{Cost:} $O(C(\mathcal{E}_n))$ barrier gradient evaluations.
\end{mdframed}

\vspace{6pt}

\begin{mdframed}[linecolor=violet!70!black, linewidth=1pt,
                 frametitle={\small\sffamily\bfseries
                 Box 3.6.3: \textsc{OnlineCoeffUpdate}}]
\label{box:online}
\textbf{Input:} $\PhaseState_n$, observable $y_n$,
cached Jacobian $J_{n-1}$, target $y_n^\star$\\[3pt]
\textbf{Step 1 - Incremental Jacobian (rank-1 secant update).}
\begin{equation}
  \tilde{J}_n
  = \frac{(y_n - y_{n-1})(\zeta_n - \zeta_{n-1})^\top}
         {\|\zeta_n - \zeta_{n-1}\|^2 + \varepsilon},
  \quad
  J_n = \rho J_{n-1} + (1-\rho)\tilde{J}_n.
  \label{eq:ch3:jacobian}
\end{equation}
The smoothed Jacobian $J_n \in \mathbb{R}^{3\times(2+|\mathcal{C}_n|)}$
is \textbf{cached} between steps and updated with a single rank-1
operation with no full recomputation.\\[3pt]
\textbf{Step 2 - Tikhonov-regularized coefficient update.}
\begin{equation}
  \Delta\zeta_n
  = (J_n^\top J_n + \lambda_J I)^{-1}
    J_n^\top\underbrace{(y_n^\star - y_n)}_{\Delta y^{\mathrm{des}}_n},
  \quad
  \zeta_{n+1}
  = \Pi_{\mathbb{R}_+}\bigl(\zeta_n + \kappa\,\Delta\zeta_n\bigr).
  \label{eq:ch3:coeff_update}
\end{equation}
The system $(J_n^\top J_n + \lambda_J I)$ is
$(2+|\mathcal{C}_n|)\times(2+|\mathcal{C}_n|)$ typically $5$--$15$
dimensional so the solve costs $O(|\mathcal{C}_n|^3)$, not $O(|\xi|^3)$.\\[3pt]
\textbf{Step 3 - Port correction (if residual $r_n$ is large).}
\begin{equation}
  r_n = \Delta y^{\mathrm{des}}_n - J_n\Delta\zeta_n,
  \qquad
  u_{f,n}^\xi
  = (P_n^\top P_n + \lambda_u I)^{-1} P_n^\top r_n.
  \label{eq:ch3:port_update}
\end{equation}
\textbf{Output:} updated $\zeta_{n+1}$, cached $J_n$, port term $u_{f,n}^\xi$\\
\textbf{What is cached:} $J_n$ (smoothed Jacobian), $\zeta_n$, $y_n$.\\
\textbf{What is wiped between stages:} inactive $\alpha_i$ (not in $\mathcal{C}_n$).\\
\textbf{What is never recomputed from scratch:} $J_n$ is always a rank-1
update of $J_{n-1}$; neither $g_\xi$ nor the Hamiltonian energy terms
are reoptimized per step.
\end{mdframed}

\begin{figure}[htbp]
    \centering
    \includegraphics[width=\linewidth]{figs/incremental_hamiltonian_update_vs_replanning.png}
    \caption{
    % \textbf{Incremental Hamiltonian update versus from-scratch replanning.}
    \textit{Left:} newly observed obstacles trigger repeated global replanning under a fixed evaluator.
    \textit{Right:} GRL-SNAM updates only the active local Hamiltonian structure, reusing cached coefficients and context to produce the phase-space transition \((q_t,p_t)\mapsto(q_{t+1},p_{t+1})\) without global recomputation.
    }
    \label{fig:ch3:incremental_vs_replan}
\end{figure}

% =======================================================================
\section{Training Losses and What Each Shapes}
\label{sec:ch3:losses}
% =======================================================================

The meta-navigator $g_\xi$ is trained offline to minimize a four-term loss.
Each term maps directly to one Hamiltonian component and one force channel;
the connection is made explicit below.

\paragraph{$\mathcal{L}_{\mathrm{traj}}$: trajectory supervision.}
$\mathcal{L}_{\mathrm{traj}} = \mathbb{E}_\mathcal{E}[\sum_t \|\vec{q}_t^{\mathrm{pred}} - \vec{q}_t^{\mathrm{ref}}\|^2]$ shapes the goal-attraction
basin ($F_{\mathrm{goal}} = -2\beta(\vec{c}-\vec{c}_{\mathrm{target}})$)
and the deformation channel ($F_{\mathrm{obj}} = -\lambda\nabla E_{\mathrm{obj}}$) via gradients that flow through the rollout
sensitivity recursion into $\partial\Potential/\partial\beta$
and $\partial\Potential/\partial\lambda$. The loss acts on phase-space trajectories $\vec{q}_t$ via the rollout
integrator, not on raw actions.
Differentiating through the symplectic integrator ensures that the
gradient shapes the \emph{force field} rather than only the action
distribution at individual timesteps which is exactly the distinction
between field learning and policy regression.

\paragraph{$\mathcal{L}_{\mathrm{vel}}$: velocity supervision.}
$\mathcal{L}_{\mathrm{vel}} = \mathbb{E}_\mathcal{E}[\sum_t \|\Mass^{-1}\Mom_t - \vec{v}_t^{\mathrm{ref}}\|^2]$
shapes the kinetic term via $\partial\Mom_t/\partial\xi$.
Without velocity supervision, trajectory loss can be minimized by
high-frequency oscillations that match positions but not momenta which is an
aliasing artifact of discrete integration.

\paragraph{$\mathcal{L}_{\mathrm{fric}}$: passivity enforcement.}
$\mathcal{L}_{\mathrm{fric}} = w_\gamma|\gamma - \gamma_0|^2$
anchors the dissipation coefficient $\gamma$ to a reference $\gamma_0$ derived from simulation.
It directly enforces the passivity condition
$\dot{H} \leq \Mom^\top\Mass^{-1} u_f$
(Definition~\ref{def:passivity}): if $\gamma$ is too small,
the system is under-damped and $\Delta H_t > 0$ along rollouts
(energy increases, passivity violated); if $\gamma$ is too large,
the system is over-damped and makes no forward progress.
$\mathcal{L}_{\mathrm{fric}}$ is \textbf{not} an arbitrary quadratic,
it is the training-time enforcement of port-Hamiltonian passivity.
$\gamma_0$ is estimated from noisy simulation rollouts.
An unweighted $|\gamma - \gamma_0|^2$ treats all estimates equally;
a covariance-weighted version $(\gamma - \gamma_0)^\top\Sigma_\gamma^{-1}(\gamma-\gamma_0)$
down-weights high-variance estimates.
The current implementation uses fixed $w_\gamma$; moving to adaptive
weighting is identified as a near-term extension in §\ref{sec:ch7:limits}.

\paragraph{$\mathcal{L}_{\mathrm{multi}}$: near-contact robustness.}
$\mathcal{L}_{\mathrm{multi}} = \frac{1}{M}\sum_m b(r_{\min} - \min_{i,t}d_i(\tilde{q}_t^{(m)};\mathcal{E}))$
shapes $\{\alpha_i\}$ and the
gradient of the IPC potential near $\{d_i = 0\}$ by explicitly over-sampling near-contact
configurations.
Without it, the learned barrier is too soft near collision because
the trajectories that matter most for collision safety are those
starting near obstacles, which are under-represented in
nominal rollouts.
$\mathcal{L}_{\mathrm{multi}}$ explicitly over-samples these
near-contact configurations.

\subsection{Full Training Objective}
\label{subsec:ch3:full_loss}

\begin{equation}
  \mathcal{L}(\xi)
  = w_t\,\mathcal{L}_{\mathrm{traj}}
  + w_v\,\mathcal{L}_{\mathrm{vel}}
  + w_\gamma\,\mathcal{L}_{\mathrm{fric}}
  + w_m\,\mathcal{L}_{\mathrm{multi}},
  \label{eq:ch3:full_loss}
\end{equation}
with the term-to-Hamiltonian mapping:
$\mathcal{L}_{\mathrm{traj}} \to \{\beta,\lambda\}$;\;
$\mathcal{L}_{\mathrm{vel}} \to \Mass^{-1}$;\;
$\mathcal{L}_{\mathrm{fric}} \to \gamma$;\;
$\mathcal{L}_{\mathrm{multi}} \to \{\alpha_i\}$.

\begin{figure}[htbp]
    \centering
    \includegraphics[width=\linewidth]{figs/loss_to_hamiltonian_structure.png}
    \caption{
    Each training loss shapes a specific Hamiltonian component, which induces a corresponding force channel and observable behavior in the geometry-only learner.
    }
    \label{fig:ch3:loss_to_hamiltonian}
\end{figure}

% =======================================================================
\section{Incremental Updates, Cache Management, and Efficiency}
\label{sec:ch3:efficiency}
% =======================================================================

A key computational property of GRL-SNAM is that \emph{nothing
is recomputed from scratch at inference time} as illustrated in Figure \ref{fig:ch3:incremental_vs_replan}.
Table~\ref{tab:ch3:cache} makes the memory and update schedule explicit.

\begin{table}[htbp]
\centering
\caption{%
  \textbf{What is stored, updated incrementally, and wiped
  in GRL-SNAM.}
  The per-step cost is dominated by the $O(|\mathcal{C}_n|^3)$
  Tikhonov solve, not by any full re-optimization.
}
\label{tab:ch3:cache}
\small
\setlength{\tabcolsep}{5pt}
\renewcommand{\arraystretch}{1.3}
\resizebox{\textwidth}{!}{
\begin{tabular}{llll}
\toprule
\textbf{Object} & \textbf{Stored where} & \textbf{Updated how} &
\textbf{Cost per step} \\
\midrule
$\xi$ (meta-navigator weights) & Offline, fixed & Not updated at inference & — \\
$\mathcal{E}_\tau$ (context) & Per-stage cache & Full recompute at $T_y$ boundary &
  $O(C)$ distance evals \\
$\zeta_n = (\eta_n,\mu_n)$ (coefficients) & Per-step cache &
  Rank-1 Tikhonov solve \eqref{eq:ch3:coeff_update} &
  $O(|\mathcal{C}_n|^3)$ \\
$J_n$ (Jacobian estimate) & Per-step cache &
  Rank-1 secant update \eqref{eq:ch3:jacobian} &
  $O(3\,|\mathcal{C}_n|)$ \\
$(\vec{q}_n,\Mom_n)$ (phase state) & Per-step &
  Leapfrog~\eqref{eq:ch3:half_mom}--\eqref{eq:ch3:full_mom} &
  $O(|\mathcal{C}_n|)$ grad evals \\
Inactive $\alpha_i$ ($i \notin \mathcal{C}_n$) & Wiped &
  Zeroed when obstacle leaves window & $O(1)$ \\
\bottomrule
\end{tabular}}
\end{table}

\paragraph{Why this is efficient.}
The meta-navigator $g_\xi$ runs once per stage as a forward pass:
$O(1)$ in the number of integration steps.
The Jacobian $J_n$ is a rank-1 update of $J_{n-1}$:
no full Jacobian is ever computed; only the direction
$(y_n - y_{n-1})(\zeta_n-\zeta_{n-1})^\top/\|\cdot\|^2$ is formed.
The coefficient solve is a $(2+|\mathcal{C}_n|)$-dimensional
Tikhonov problem, not a full network optimization.
In a typical corridor, $|\mathcal{C}_n| \leq 8$, making the
per-step cost trivially small.

\paragraph{Why this is more robust than fixed estimation.}
The smoothed Jacobian $J_n = \rho J_{n-1} + (1-\rho)\tilde{J}_n$
accumulates a running estimate of how the observable $y_t$ responds
to coefficient changes.
This is a stateful, incrementally improving estimator and not a
Jacobian computed once at initialization and held fixed.
As the rollout progresses and more $(y,\zeta)$ pairs are observed,
$J_n$ becomes more accurate, and the Tikhonov updates become more
targeted.
This is the sense in which GRL-SNAM provides
\emph{stateful estimator improvement} without storing full trajectories.

% =======================================================================
\section{Strengths of the Geometry-Only Learner}
\label{sec:ch3:strengths}
% =======================================================================

\paragraph{Explicit field topology.}
The factored potential \eqref{eq:ch3:potential} makes the force-field
topology inspectable: goal-attracting basins, repulsive walls, and
dissipative channels are separate terms with separate coefficients.
Any coefficient can be examined or corrected without touching the others.

\paragraph{Stable online correction.}
The passivity condition $\dot{H} \leq \Mom^\top\Mass^{-1}u_f$
is enforced at training time by $\mathcal{L}_{\mathrm{fric}}$
and monitored at inference by $\mathcal{L}_{\mathrm{pass}}$.
Online updates \eqref{eq:ch3:coeff_update} respect the
$\Pi_{\mathbb{R}_+}$ projection, keeping all coefficients
non-negative and preserving the repulsive character of each barrier.

\paragraph{Sample efficiency.}
500k gradient steps of Hamiltonian-structured supervision achieve
87.5\% success on the dungeon benchmark which is an order of magnitude
fewer than the 3--8M steps required by PPO/SAC/TRPO baselines
under the same interface.

\paragraph{Minimal sensing effort.}
$\rho_{\mathrm{map}} \leq 0.35$ on all test splits: the agent
explores less than 35\% of the workspace because the energy landscape
is refined only in the neighborhood actually visited.

% =======================================================================
\section{Limitations of Geometry-Only Structure}
\label{sec:ch3:limitations}
% =======================================================================

Three structural gaps follow directly from the geometry-only
restriction~\eqref{eq:ch3:context}.
These are not tuning failures, they are architectural.

\paragraph{L1: No material-risk force channel.}
The potential \eqref{eq:ch3:potential} has no term for $\tilde{r}(\vec{c})$.
Two regions with identical clearances but different material risk
(safe grass vs.\ hazardous mud) receive identical forces.
Adding a material-risk channel requires a new energy term, a new
coefficient, a new gradient pathway, and a new loss term which is 
precisely the enrichment of Chapter~\ref{chap:material}.

\paragraph{L2: No tail-risk objective.}
The training loss \eqref{eq:ch3:full_loss} optimizes expected
trajectory quality.
It cannot focus gradient signal on rare catastrophic outcomes
(the tails of the rollout-cost distribution).
This requires CVaR, a coherent risk measure not expressible as
an expectation, together with the material-aware episode cost $J$.

\paragraph{L3: Coefficient space saturated.}
The coefficient vector $(\beta,\lambda,\{\alpha_i\},\gamma)$ controls
only the weights of existing geometric terms.
Adding a new \emph{type} of force (material repulsion, hazard barrier)
requires a new slot in this vector, a new head in $g_\xi$, and a new
input to the online update.
None of this fits in the current coefficient space.

\begin{quote}\itshape
Geometry-only structured Hamiltonian learning yields a principled,
sample-efficient, continually correctable navigation policy under
local sensing but the resulting learner is blind to latent semantic
risk, lacks a tail-risk objective, and has no energy-coefficient
slot for material information.
Chapter~\ref{chap:material} addresses all three by adding new energy
terms, following the enrichment chain of Theorem~\ref{thm:ch5:enrichment}.
\end{quote}